% Case studies. Every number is a macro from analysis/p7_cases.py, which
% reads the medians of the measured runs and asserts full separation.
\paragraph{Faster, but not for the claimed reason}
iPhotron PR~\#453 (Copilot) claims to remove the slowdown of deep SQL
pagination and to load assets on demand instead of all at once. The workload
runs \caseAWall\ of the base wall time (\caseAWallBase\,s to
\caseAWallHead\,s; all ratios use unrounded medians), yet the kernel shows more work, not less: \caseASyscalls\
the syscalls (\caseASyscallsBase\ to \caseASyscallsHead), \caseAPread\ the
\texttt{pread64} calls, \caseABytes\ the bytes read and \caseASys\ the system
time. Total CPU time grows \caseATask\ while the process runs
\caseAThreadsHead\ threads instead of \caseAThreadsBase. Peak memory does not
change. A stopwatch would credit the claim; the trace does not support it.

\paragraph{Slower, and the trace shows where}
iPhotron PR~\#564 (Jules) uses pre-generated thumbnails. On our workload,
which keeps 500 unrelated files in the thumbnail folder, it is
\caseBWall\ slower (\caseBWallBase\,s to \caseBWallHead\,s) and its output
changes. The trace shows a likely cause: \texttt{getdents64}, the call that
lists a directory, rises from \caseBGetdentsBase\ to \caseBGetdentsHead\
(\caseBGetdents), and instructions grow \caseBInstr.

\paragraph{Half the mechanism happens, and it hurts}
A Claude Code PR to gemini-notebook-mcp-cli (\#16) claims non-blocking
I/O and HTTP connection pooling. The first part appears: the main thread
spends \caseDOffcpu\ of its former time off the CPU. The second does not:
the workload opens \caseDSocketHead\ sockets and \caseDConnectHead\
connections on both versions. Meanwhile syscalls grow \caseDSyscalls,
\texttt{futex} calls \caseDFutexc\ and CPU time \caseDTask, and the
workload is \caseDWall\ slower (\caseDWallBase\,s to \caseDWallHead\,s).

\paragraph{A claim that holds}
mlflow PR~\#19520 (Copilot) changes artifact transfer to chunked streaming.
The footprint is exactly the claimed one: \texttt{read} calls fall from
\caseCReadBase\ to \caseCReadHead\ (\caseCRead) for the same bytes read,
instructions fall to \caseCInstr\ and wall time to \caseCWall\ of the base.
